\section{Related Work}
\label{sec:related}

\subsection{Honeypot and Honeynet Architectures}
Deception research has moved from single-service traps to orchestrated
honeynets~\cite{spitzner2002}, with recent surveys cataloging managed
deception fabrics and cloud-native deployment~\cite{lopez2024}, from
multi-sensor platforms that aggregate dozens of low- and
medium-interaction sensors in one stack to containerized honeynets that
scale the same idea elastically. The survey's gap analysis is blunt: no
widely deployed framework grounds its responses in real-time threat
intelligence or differentiates behavior across attacker classes, which
are the two capabilities G1 and \S\ref{sec:chrollo} exist to provide.
What these architectures share is a static
relationship between the decoy and its content: the deception is authored
once, at deployment time. \system differs in kind, not degree: its served
content is decided per session, and every decision passes a verification
gate before reaching the wire.

\subsection{Generative Models in Deception}
LLM-backed honeypots are the closest line of work. GPT-driven SSH bots
demonstrate that language models can hold convincing interactive
sessions~\cite{valeros2024}; recent systems couple generation to
retrieval for grounded responses~\cite{qiao2025,liang2025}, and agentic
RAG pipelines have been proposed for
CTI workflows~\cite{belmahjoub2025,gao2025}. This literature establishes
that LLM honeypots are feasible and that grounding helps. What it does not
establish is \emph{what happens when the generator is wrong}: the
consistency, safety, and degradation questions that a security tool must
answer before such a system can be left facing an attacker. Those
questions are this paper's subject, and the fiction graph, serving-time
verification, and deterministic fallback are our answers. LLMPot
extends LLM emulation to industrial protocols with dynamic
configuration~\cite{vasilatos2025llmpot}, strengthening the case for
generative deception while leaving the serving-trust question open;
independent honeypots whose state must remain consistent across
distributed layers reach for explicit rule machinery of their
own~\cite{guan2025cpds}. Recent runtime-enforcement work places the
boundary similarly: the probabilistic component proposes, while an
external, inspectable mechanism owns the safety
decision~\cite{wang2026agentspec} --- \system applies that principle to
served generative deception itself, where the verified object is the
attacker-facing response rather than an agent action or a
protocol state. The
methodological difference is as important as the architectural one:
Kouremetis et al.~\cite{kouremetis2024} and Qiao et al.~\cite{qiao2025}
report prompt-engineering improvements, whereas
our evaluation isolates each mechanism with paired ablations and reports
the ablation that correctly produced a null result (\S\ref{sec:rq1}).

\subsection{LLM Security and Alignment Pressure}
Our provider-fragility observations connect to a broader alignment literature.
Safety-aligned models refuse or deflect on adversarial-role prompts, and
constitutional and RLHF training shapes exactly the refusal behavior our
harness must route around~\cite{bai2022,ganguli2022,christiano2017}.
Prompt-injection risks in retrieval pipelines motivate the corpus hygiene
of \S\ref{sec:implementation}~\cite{greshake2023,perez2022}; structured-query
and preference-optimization defenses harden the model against
injected instructions~\cite{chen2025struq,chen2025secalign}, and
retrieval-poisoning studies show that a handful of poisoned documents
in a large corpus dominate retrieval
rankings~\cite{zou2025poisonedrag,malladi2026ripeii}, with
recent work beginning to attribute misleading generations to the
responsible corpus evidence~\cite{zhang2026ragorigin}. Our curated,
provenance-screened corpus and
channel separation treat retrieved evidence as untrusted input, and
multi-tool agent studies extend the attack surface from model output to
orchestration state~\cite{li2026dissonances}. The step this paper adds
is architectural: even a robust generator's output does not reach the
wire before deterministic verification, so the boundary does not depend
on robustness holding~\cite{chen2025struq,chen2025secalign}. And
documented failures of ML-backed security tooling under distribution
shift~\cite{reworr2024} reinforce the case for measured fallbacks over
trusting generation. The deployment lesson we add is operational:
treat the provider as an unreliable component with catalog drift, quota
exhaustion, and alignment-induced timeouts as normal operating conditions.
The reviewer-facing implication is a methodological one: evaluations of
LLM-driven defenses should report provider instability as part of the
results, because it shaped which experiments were even runnable.

\subsection{CTI and Behavior Analysis}
Hybrid RAG over threat intelligence has been shown effective for
attribution and technique mapping~\cite{alsada2024,rajapaksha2024},
and behavioral profiling of attacker sessions underpins
honeypot-facing classification~\cite{lekssays2025}. Adaptive
deception experiments show that attacker skill modulates effective
strategy~\cite{kouremetis2024}, and honeytoken-based
attribution motivates the canary-registered bait
design. \system composes these primitives into a loop whose safety
properties, rather than its raw capability, are the contribution.
